
In dieser Arbeit wird ein Einspurmodell zur Beschreibung der lateralen Fahrzeugdynamik
untersucht. Grundlage der Analyse sind sowohl ein nichtlineares als auch ein linearisiertes
Modell, welche in Simulink umgesetzt wurden.

Der Arbeit liegen zwei Simulink-Dateien bei. Die erste Datei enthält sowohl das
nichtlineare als auch das linearisierte Einspurmodell. Diese wurde für die Aufgaben~2 und~3
verwendet, in denen die Modellgleichungen implementiert sowie das Verhalten des linearen
und des nichtlinearen Systems miteinander verglichen wird.

Die zweite Simulink-Datei enthält ausschließlich das linearisierte Modell und wurde für die
Untersuchung des Frequenzverhaltens in den Aufgaben~4 und~5 verwendet, insbesondere für die
Bestimmung des Frequenzgangs und die Darstellung der Bodediagramme.

Zusätzlich ist ein MATLAB-Live-Skript beigefügt, in dem die Aufgaben~1, 4 und~5 bearbeitet
wurden. Darin sind die analytischen Berechnungen sowie die Auswertung der Simulationen
und Frequenzgänge umgesetzt.

Die Grundlage des verwendeten Modells bilden die folgenden nichtlinearen Differentialgleichungen
des Einspurmodells, wie sie in der Aufgabenstellung gegeben sind. Diese beschreiben die
Schräglaufwinkel der Reifen sowie die Dynamik des Schwimmwinkels $\beta$ und der Gierrate
$\dot{\Psi}$.

\subsection*{Nichtlineare Modellgleichungen}

Dynamik des Schwimmwinkels:
\begin{equation}
\dot{\beta}
=
-\dot{\Psi}
+
\frac{1}{m v}
\left(
c_v \alpha_v \frac{\cos(\delta)}{\cos(\beta)}
+
c_h \alpha_h \frac{1}{\cos(\beta)}
\right)
\label{eq:betadot}
\end{equation}
\\
Dynamik der Gierrate:
\begin{equation}
\ddot{\Psi}
=
\frac{1}{J}
\left(
l_v c_v \alpha_v \cos(\delta)
-
l_h c_h \alpha_h
\right)
\label{eq:psiddot}
\end{equation}
\\
Schräglaufwinkel an der Vorderachse:
\begin{equation}
\alpha_v
=
\delta
-
\frac{l_v \dot{\Psi}}{v \cos(\beta)}
+
\tan(\beta)
\label{eq:alphav}
\end{equation}
\\
Schräglaufwinkel an der Hinterachse:
\begin{equation}
\alpha_h
=
\frac{l_h \dot{\Psi}}{v \cos(\beta)}
-
\tan(\beta)
\label{eq:alphah}
\end{equation}
\\
Dabei bezeichnet $m$ die Fahrzeugmasse, $v$ die Fahrzeuggeschwindigkeit,
$J$ das Trägheitsmoment um die Hochachse, $l_v$ und $l_h$ die Abstände des
Schwerpunkts zu Vorder- bzw. Hinterachse sowie $c_v$ und $c_h$ die
Reifensteifigkeiten der Vorder- und Hinterachse.